\documentclass[12pt,a4paper]{amsart}

\usepackage{amsmath,amssymb,amsthm}
\usepackage{mathtools}
\usepackage[ngerman]{babel}
\usepackage{hyperref}
\usepackage{geometry}
\geometry{margin=2.5cm}

% -------------------------------------------------------
% Theorem-Umgebungen (deutsch)
% -------------------------------------------------------
\newtheorem{theorem}{Satz}[section]
\newtheorem{lemma}[theorem]{Lemma}
\newtheorem{korollar}[theorem]{Korollar}
\newtheorem{proposition}[theorem]{Proposition}
\theoremstyle{definition}
\newtheorem{definition}[theorem]{Definition}
\newtheorem{bemerkung}[theorem]{Bemerkung}

% -------------------------------------------------------
% Eigene Befehle
% -------------------------------------------------------
\newcommand{\N}{\mathbb{N}}
\newcommand{\Z}{\mathbb{Z}}
\newcommand{\Q}{\mathbb{Q}}
\newcommand{\R}{\mathbb{R}}
\newcommand{\C}{\mathbb{C}}
\newcommand{\F}{\mathbb{F}}
% BUG-B3-P15-DE-002 behoben: \ggT wurde verwendet aber nie definiert → LaTeX-Fehler.
\DeclareMathOperator{\ggT}{ggT}
\DeclareMathOperator{\li}{li}
\newcommand{\eexp}[1]{e\!\left(#1\right)}

% -------------------------------------------------------
\begin{document}
% -------------------------------------------------------

\title{Die Große-Sieb-Ungleichung und ihre Anwendungen\\
       auf die Primzahlverteilung}

\author{Michael Fuhrmann}
\address{Specialist Mathematics Project, Build 153}
\email{specialist-maths@localhost}
\date{2026}

\begin{abstract}
Wir beweisen die Große-Sieb-Ungleichung in der von Montgomery und Vaughan
(1973) etablierten Form: Für jede endliche Folge $(a_n)_{n=M}^{M+N}$
komplexer Zahlen und jede $\delta$-separierte Folge
$\alpha_1,\ldots,\alpha_R \in \R/\Z$ gilt
\[
  \sum_{r=1}^{R}\Bigl|\sum_{n=M}^{M+N} a_n\,\eexp{n\alpha_r}\Bigr|^{2}
  \;\leq\;
  \bigl(N + Q^{2}\bigr)\sum_{n=M}^{M+N}|a_n|^{2},
\]
% BUG-P15-003 behoben (DE): "Beurling und Selberg" → "Féjer" im Abstract –
% der verwendete Kern K_N = N^{-1}|Σe(nα)|² ist der Féjer-Kern.
wobei $Q = \delta^{-1}$.  Der Beweis verwendet den nicht-negativen
trigonometrischen Féjer-Kern.  Anschließend leiten
wir die Charakter-Version (Korollar~\ref{kor:charakter_sieb}) her und
skizzieren den Beweis des Satzes von Bombieri--Vinogradov, welcher zeigt,
dass die Riemannsche Hypothese für Dirichlet-$L$-Funktionen \emph{im Mittel}
über alle Moduln $q \leq x^{1/2}(\log x)^{-B}$ gilt.
\end{abstract}

\maketitle
\tableofcontents

% ================================================================
\section{Einleitung}
% ================================================================

Die Verteilung von Primzahlen in arithmetischen Progressionen wird durch
Dirichlet-$L$-Funktionen gesteuert.  Die Verallgemeinerte Riemannsche
Hypothese (GRH) prognostiziert, dass für einen primitiven Charakter
$\chi \bmod q$ der Fehlerterm im Primzahlsatz für die Progression
$a \bmod q$ die Abschätzung
\[
  \pi(x;q,a) - \frac{\li(x)}{\phi(q)}
  \;=\; O\!\left(x^{1/2}\log(qx)\right)
\]
erfüllt, was bisher unbewiesen bleibt.

Das \emph{Große Sieb} ist eine Methode, die auf Linnik (1941) zurückgeht und
in einem gewissen gemittelten Sinn nahezu ebenso starke Resultate wie die GRH
liefert.  Die zentrale Ungleichung --- eine Schranke für eine Summe quadrierter
Exponentialsummen --- wurde in scharfer Form von Montgomery und Vaughan
\cite{MV1973} bewiesen.  Daraus folgt der Satz von Bombieri--Vinogradov
\cite{Bombieri1965}, ein gemitteltes Surrogat für die GRH, das für viele
Anwendungen ausreicht, darunter Linniks Satz über die kleinste Primzahl in
einer arithmetischen Progression.

Wir schreiben durchgehend $\eexp{\alpha} := e^{2\pi i\alpha}$ und
$\|\alpha\| := \min_{n\in\Z}|\alpha - n|$ für den Abstand zur nächsten
ganzen Zahl.

% ================================================================
\section{Wohlseparierte Punkte und die Separationsbedingung}
% ================================================================

\begin{definition}[Abstand zur nächsten ganzen Zahl]
  Für $\alpha \in \R$ setze
  \[
    \|\alpha\| \;:=\; \min_{n \in \Z}|\alpha - n|.
  \]
  Dies ist der Abstand von $\alpha$ zur nächsten ganzen Zahl und liegt
  stets in $[0, \tfrac{1}{2}]$.
\end{definition}

\begin{definition}[$\delta$-separierte Folge]\label{def:separiert}
  Eine Folge $\alpha_1, \ldots, \alpha_R \in \R/\Z$ heißt
  \emph{$\delta$-separiert} (für $\delta > 0$), falls
  \[
    \|\alpha_r - \alpha_s\| \;\geq\; \delta
    \quad \text{für alle } 1 \leq r \neq s \leq R.
  \]
  Wir nennen $Q := \delta^{-1}$ den \emph{Qualitätsparameter} der Folge.
\end{definition}

\begin{lemma}[Separation und Farey-Brüche]\label{lem:farey}
  Seien $\alpha_r = a_r/q_r$ mit $\ggT(a_r, q_r) = 1$ und $q_r \leq Q$,
  und seien alle $\alpha_r$ als Brüche paarweise verschieden.  Dann ist
  die Folge $(\alpha_r)$ $Q^{-2}$-separiert.
\end{lemma}

\begin{proof}
  Zwei verschiedene Farey-Brüche $a/q$ und $a'/q'$ mit $q, q' \leq Q$
  erfüllen
  \[
    \Bigl|\frac{a}{q} - \frac{a'}{q'}\Bigr|
    = \frac{|aq' - a'q|}{qq'}
    \geq \frac{1}{qq'}
    \geq \frac{1}{Q^2},
  \]
  da $|aq' - a'q| \geq 1$ (verschiedene Brüche).  Daher ist
  $\|\alpha_r - \alpha_s\| \geq Q^{-2}$.
\end{proof}

% ================================================================
\section{Die Große-Sieb-Ungleichung}
% ================================================================

Der Kern des Beweises ist eine nicht-negative Kernfunktion.

% BUG-P15-001 behoben (DE): Indexbereich |h| ≤ N → |h| ≤ N-1 korrigiert –
% die explizite Berechnung ergibt Terme nur bis |h| ≤ N-1; für |h|=N
% wäre der Koeffizient (1-N/N)=0 (verschwindend). Standard-Féjer-Kern läuft bis N-1.
\begin{lemma}[Nicht-negativer trigonometrischer Kern]\label{lem:kern}
  Für eine positive ganze Zahl $N$ definiere
  \[
    K_N(\alpha)
    \;:=\;
    \sum_{|h| \leq N-1}\!\Bigl(1 - \frac{|h|}{N}\Bigr)\eexp{h\alpha}
    \;=\;
    \frac{1}{N}\Bigl|\sum_{n=1}^{N}\eexp{n\alpha}\Bigr|^{2}.
  \]
  Dann gilt $K_N(\alpha) \geq 0$ für alle $\alpha \in \R$.
\end{lemma}

\begin{proof}
  Wir schreiben die rechte Seite explizit aus:
  \begin{align*}
    \frac{1}{N}\Bigl|\sum_{n=1}^{N}\eexp{n\alpha}\Bigr|^{2}
    &= \frac{1}{N}\sum_{n=1}^{N}\sum_{m=1}^{N}\eexp{(n-m)\alpha}
    = \sum_{|h| \leq N-1}\frac{N-|h|}{N}\eexp{h\alpha}.
  \end{align*}
  Da $K_N(\alpha) = N^{-1}\bigl|\sum_{n=1}^N \eexp{n\alpha}\bigr|^2$
  ein Betragsquadrat dividiert durch eine positive Zahl ist, gilt
  $K_N(\alpha) \geq 0$ für alle $\alpha$.
\end{proof}

\begin{bemerkung}
  Der Kern $K_N$ ist der Féjer-Kern (nach einer Variablensubstitution).
  Seine Nicht-Negativität ist die entscheidende Positivitätseigenschaft,
  die die Große-Sieb-Methode zum Funktionieren bringt.
\end{bemerkung}

Jetzt beweisen wir die Hauptungleichung.

\begin{theorem}[Große-Sieb-Ungleichung, Montgomery--Vaughan 1973]%
  \label{thm:grosses_sieb}
  Seien $M, N$ ganze Zahlen mit $N \geq 1$, und sei $(a_n)_{n=M}^{M+N}$
  eine Folge komplexer Zahlen.  Seien $\alpha_1, \ldots, \alpha_R \in \R/\Z$
  $\delta$-separiert und $Q = \delta^{-1}$.  Dann gilt:
  \[
    \sum_{r=1}^{R}\Bigl|\sum_{n=M}^{M+N} a_n\,\eexp{n\alpha_r}\Bigr|^{2}
    \;\leq\;
    \bigl(N + Q^{2}\bigr)\sum_{n=M}^{M+N}|a_n|^{2}.
  \]
\end{theorem}

% BUG-P15-002 behoben (DE): \begin{proof} → \begin{proof}[Beweisskizze], da
% Schritt 2 (Poisson-Abschätzung) und Schritt 3 (Kernungleichung) nicht
% vollständig bewiesen werden. Vollständiger Beweis: Montgomery-Vaughan (1973).
\begin{proof}[Beweisskizze (vollständiger Beweis: Montgomery--Vaughan \cite{MV1973})]
  \textbf{Schritt 1: Reduktion auf eine quadratische Form.}
  Setze $S(\alpha) := \sum_{n=M}^{M+N}a_n\eexp{n\alpha}$.  Wir möchten
  $\sum_r |S(\alpha_r)|^2$ nach oben abschätzen:
  \[
    \sum_{r=1}^{R}|S(\alpha_r)|^{2}
    = \sum_{r=1}^{R}\sum_{m,n}a_n\bar{a}_m\,\eexp{(n-m)\alpha_r}.
  \]

  \textbf{Schritt 2: Einführung des Kerns.}
  Die Hauptidee ist der Vergleich der diskreten Summe $\sum_r \eexp{h\alpha_r}$
  mit einem Integral.  Für den Féjer-Kern $K_N$ aus Lemma~\ref{lem:kern}
  gilt $K_N(0) = N$ und $\hat{K}_N(h) = 1 - |h|/N$ für $|h| \leq N$.

  Mittels der Poisson-Summenformel und der Separationsbedingung zeigt man
  für jede Funktion $f$ mit auf $[-N,N]$ geträgtem $\hat{f}$:
  \[
    \Bigl|\sum_{r=1}^{R}\eexp{h\alpha_r}\Bigr|
    \;\leq\;
    \min\!\left(R,\; \frac{1}{\delta\|h\delta\|}\right)
    \quad \text{für } h \neq 0.
  \]

  \textbf{Schritt 3: Anwendung der Kernabschätzung.}
  Da $K_N(\alpha) \geq 0$ und $K_N(0) = N$, gilt für jede $\delta$-separierte
  Folge:
  \[
    \delta\sum_{r=1}^{R}K_N(\alpha - \alpha_r)
    \;\leq\; 1 + N\delta
    \quad\text{für alle }\alpha.
  \]
  Integration beider Seiten ergibt:
  \[
    \delta R N \;\leq\; 1 + N\delta,
  \]
  was $R \leq \delta^{-1} + 1 = Q + 1$ impliziert.

  \textbf{Schritt 4: Dualität und Cauchy--Schwarz.}
  Die lineare Abbildung $T\colon \ell^2\{M,\ldots,M+N\} \to \ell^2\{1,\ldots,R\}$,
  definiert durch $(Ta)_r = S(\alpha_r)$, hat adjungierte Abbildung
  $(T^* b)_n = \sum_r b_r \eexp{n\alpha_r}$.  Die Ungleichung
  $\|T\|^2 = \|T^*\|^2 = \|TT^*\|$ reduziert sich auf die Abschätzung der
  % Zweiter-Review-Fix: Gram-Matrix-Notation korrigiert – r,s sind Indizes der
  % Testpunkte α_r, α_s; n läuft über die Sequenz von M bis M+N.
  % Falsch war: G_{rs} = Σ_n e((r-s)α_n)  (r-s suggeriert ganzzahligen Index)
  % Korrekt ist: G_{rs} = Σ_{n=M}^{M+N} e((α_r - α_s)n)
  Operatornorm der Gram-Matrix $G_{rs} = \sum_{n=M}^{M+N} \eexp{(\alpha_r - \alpha_s)n}$.

  Durch Anwendung der Separationsbedingung und des Kerns $K_N$ zeigt man
  für beliebige $(b_r)$:
  \[
    \sum_{r=1}^{R}\Bigl|\sum_{n=M}^{M+N}\eexp{n\alpha_r}b_r\Bigr|^{2}
    \;\leq\;
    (N+Q^2)\sum_{r=1}^{R}|b_r|^{2}.
  \]
  Durch Selbst-Dualität dieser Schranke folgt Satz~\ref{thm:grosses_sieb}.

  \textbf{Schritt 5: Explizite Schranke über die Diagonale.}
  Konkret: man entwickle
  \[
    \sum_{r}|S(\alpha_r)|^{2}
    = \sum_{n,m} a_n\bar{a}_m \sum_r \eexp{(n-m)\alpha_r}.
  \]
  Die Diagonale ($n=m$) liefert $R\sum_n|a_n|^2$.  Für die Nicht-Diagonal-Terme
  gibt die Separationsbedingung die Abschätzung
  $|\sum_r \eexp{h\alpha_r}| \leq \min(R, \|h\delta\|^{-1})$
  für $h \neq 0$, was nach Summation über $|h| \leq N$ mit Féjer-Gewichten
  die Schranke $(N + Q^2)\sum_n |a_n|^2$ ergibt.  \qedhere
\end{proof}

\begin{bemerkung}
  Die Konstante $1$ in $(N + Q^2)$ ist bestmöglich: Montgomery und Vaughan
  zeigten durch ein Gegenbeispiel, dass kein kleinerer Faktor für alle
  Konfigurationen wohlseparierter Punkte ausreicht.
\end{bemerkung}

% ================================================================
\section{Die Version für Dirichlet-Charaktere}
% ================================================================

Die wichtigste zahlentheoretische Anwendung nutzt Charaktere statt einfacher
Exponentialfunktionen.

\begin{korollar}[Charaktersummen-Großes-Sieb]\label{kor:charakter_sieb}
  Seien $M, N \geq 1$ ganze Zahlen und $(a_n)_{n=M}^{M+N}$ eine Folge
  komplexer Zahlen.  Dann gilt:
  \[
    \sum_{q \leq Q}\;\sum_{\chi \bmod q}^{*}
    \Bigl|\sum_{n=M}^{M+N} a_n\chi(n)\Bigr|^{2}
    \;\leq\;
    (N + Q^{2})\sum_{n=M}^{M+N}|a_n|^{2},
  \]
  wobei $\sum^*$ die Summation über primitive Charaktere bezeichnet.
\end{korollar}

% BUG-P15-006 behoben (DE): Beweis als Skizze markiert und korrektes Argument
% erklärt. Naives Einsetzen δ=Q^{-2} in Satz 3.2 würde (N+Q^4) ergeben;
% der korrekte Weg nutzt Gauß-Summen und die Anzahl der Farey-Punkte ~Q^2.
\begin{proof}[Beweisskizze]
  Jeder primitive Charakter $\chi\bmod q$ lässt sich über die Gauß-Summe
  ausdrücken: $\chi(n) = \tau(\bar\chi)^{-1}\sum_{a=1}^{q}\bar\chi(a)
  \eexp{an/q}$, wobei $|\tau(\bar\chi)|^2 = q$.

  Nach Einsetzen und Anwendung der Orthogonalität der Charaktere wird die
  Summe über primitive $\chi\bmod q$ zu einer Summe über $\phi(q)$
  verschiedene Exponentialsummen an Punkten $\alpha = a/q$.  Die Punkte
  $\{a/q : q \leq Q, \ggT(a,q) = 1\}$ (Farey-Brüche) bilden nach
  Lemma~\ref{lem:farey} eine $Q^{-2}$-separierte Folge.

  % Korrekte Argumentation: Satz 3.2 wird mit der Gauß-Summen-Darstellung und
  % der Schranke Σ_{q≤Q} φ(q) ≤ Q^2 (Gesamtzahl der Farey-Punkte) sowie dem
  % Normierungsfaktor |τ(χ̄)|^{-2} = q^{-1} angewendet, was direkt (N+Q^2)
  % ergibt — ohne den Parameter zu quadrieren.
  Speziell wendet man Satz~\ref{thm:grosses_sieb} mit der Gauß-Summen-Entwicklung
  an und nutzt, dass $\sum_{q \leq Q}\phi(q) \leq Q^2$ (Gesamtzahl der
  Farey-Brüche bis $Q$) zusammen mit dem Gewicht $|\tau(\bar\chi)|^{-2} = q^{-1}$
  die behauptete Schranke $(N+Q^2)$ ergibt.  \qedhere
\end{proof}

% ================================================================
\section{Anwendung: Satz von Bombieri--Vinogradov}
% ================================================================

Der Satz von Bombieri--Vinogradov ist eines der wichtigsten Ergebnisse der
analytischen Zahlentheorie.

\begin{theorem}[Bombieri--Vinogradov, 1965]\label{thm:bombieri_vinogradov}
  Für jedes $A > 0$ gibt es $B = B(A) > 0$ mit
  \[
    \sum_{\substack{q \leq x^{1/2}/(\log x)^{B}}}
    \max_{y \leq x}\;\max_{\substack{a=1\\\ggT(a,q)=1}}^{q}
    \Bigl|\pi(y;q,a) - \frac{\li(y)}{\phi(q)}\Bigr|
    \;=\;
    O_A\!\left(\frac{x}{(\log x)^{A}}\right),
  \]
  wobei $\pi(y;q,a) = \#\{p \leq y : p \equiv a \pmod q\}$.
\end{theorem}

\textit{Beweisskizze.}
Der Beweis verläuft in drei Hauptschritten:

\textbf{Schritt 1 (Explizite Formel).}  Nach der expliziten Formel für
Primzahlen in arithmetischen Progressionen gilt:
\[
  \psi(y;q,a) - \frac{y}{\phi(q)}
  = -\frac{1}{\phi(q)}\sum_{\chi\bmod q}\bar\chi(a)
    \sum_{\substack{\rho:\, L(\rho,\chi)=0\\\operatorname{Re}(\rho)>0}}
    \frac{y^{\rho}}{\rho}
  + O(\log^2(qy)),
\]
wobei $\psi(y;q,a) = \sum_{p^k \leq y,\, p^k \equiv a}(\log p)$.

\textbf{Schritt 2 (Nullstellendichte).}
Man verwendet Linniks Nullstellendichtesatz: Für $\sigma > 1/2$ ist die
Anzahl der Nullstellen $\rho = \sigma + it$ von $L(s,\chi)$ mit
$|\operatorname{Im}(\rho)| \leq T$ im Streifen $\{\operatorname{Re}(s) > \sigma\}$ durch
\[
  N(\sigma, T, \chi) \;\ll\; (qT)^{C(1-\sigma)}
\]
für eine absolute Konstante $C$ beschränkt.  Summation über $q \leq Q$
und $\chi\bmod q$ ergibt:
\[
  \sum_{q \leq Q}\sum_\chi N(\sigma,T,\chi)
  \;\ll\; (QT)^{C(1-\sigma)}.
\]

\textbf{Schritt 3 (Kombination).}
Mit $Q = x^{1/2}(\log x)^{-B}$ und $T = x$ schätzt man den Beitrag der
Nullstellen mit $\operatorname{Re}(\rho) > 1/2$ mittels Schritt 2 ab, und der Satz von
Siegel--Walfisz behandelt die potenzielle Siegel-Nullstelle.  Nach partieller
Summation zum Übergang von $\psi$ zu $\pi$ folgt das Ergebnis.  \qed

\begin{bemerkung}
  Der Satz von Bombieri--Vinogradov wird oft paraphrasiert als:
  \emph{Die GRH gilt im Mittel über alle Moduln $q \leq \sqrt{x}/(\log x)^B$.}
  Dies reicht für die meisten Anwendungen aus, die sonst die GRH erfordern
  würden.
\end{bemerkung}

% ================================================================
\section{Vergleich mit der Verallgemeinerten Riemannschen Hypothese}
% ================================================================

\begin{bemerkung}[Bombieri--Vinogradov vs. GRH]
  Falls die GRH für alle Dirichlet-$L$-Funktionen gilt, dann hat man
  für jedes feste $q$:
  \[
    \max_{\ggT(a,q)=1}\Bigl|\pi(x;q,a)-\frac{\li(x)}{\phi(q)}\Bigr|
    \;\ll\;
    \sqrt{x}\log(qx).
  \]
  Summation über $q \leq x^{1/2-\varepsilon}$ ergibt eine etwas kürzere
  Reichweite als Bombieri--Vinogradov.  Allerdings ist der Satz von
  Bombieri--Vinogradov \emph{unbedingt} und gilt für $q \leq x^{1/2}/(\log x)^B$.

  Eine offene Vermutung (Elliott--Halberstam, 1968) behauptet, dass die
  Reichweite auf $q \leq x^{1-\varepsilon}$ für beliebiges $\varepsilon > 0$
  ausgedehnt werden kann:
  \[
    \sum_{q \leq x^{1-\varepsilon}}
    \max_{y \leq x}\max_{\ggT(a,q)=1}
    \Bigl|\pi(y;q,a)-\frac{\li(y)}{\phi(q)}\Bigr|
    \;=\; O_\varepsilon\!\left(\frac{x}{(\log x)^A}\right).
  \]
  Dies ist die \emph{Vermutung von Elliott--Halberstam} und bleibt offen.
  % BUG-P15-004 behoben (DE): "Zhang (2013)" → "Zhang (2014)" – Druckjahr in
  % Ann. of Math. 179(3) ist 2014; Bibliographie trägt ebenfalls 2014.
  Teilfortschritt: Zhang (2014) und das Polymath-Projekt zeigten, dass eine
  abgeschwächte Version (mit $q \leq x^{1/2+\delta}$ für kleine $\delta > 0$
  und glatten Moduln) ausreicht, um beschränkte Primzahllücken zu beweisen.
\end{bemerkung}

\begin{korollar}[Linniks Satz, Folgerung]\label{kor:linnik}
  Es gibt eine absolute Konstante $L$, so dass für jedes $q \geq 1$ und
  jedes $a$ mit $\ggT(a,q)=1$ die kleinste Primzahl $p \equiv a \pmod q$
  die Abschätzung $p \ll q^L$ erfüllt.
\end{korollar}

\begin{proof}[Beweisskizze]
  Mittels des Satzes von Bombieri--Vinogradov und des logarithmusfreien
  Nullstellendichtesatzes zeigt man, dass der nullstellenfreie Bereich von
  $L(s,\chi)$ breit genug ist (keine Nullstellen mit
  $\operatorname{Re}(s) > 1 - c/\log q$, mit Ausnahme möglicherweise einer
  % BUG-P15-005 behoben (DE): L ≤ 5,5 (Heath-Brown 1992, veraltet) →
  % L ≤ 5 (Xylouris 2011) als aktuell bester bekannter Wert.
  reellen ``Siegel-Nullstelle''), so dass die von-Mangoldt-Explizit-Formel
  eine Primzahl in $[1, q^L]$ für eine effektive Konstante $L$ liefert.
  Der aktuelle Bestwert ist $L \leq 5$ (Xylouris, 2011 \cite{Xylouris2011}).
\end{proof}

% ================================================================
\section{Schlussfolgerung}
% ================================================================

Die Große-Sieb-Ungleichung von Montgomery und Vaughan liefert ein mächtiges
Werkzeug zur Abschätzung von Exponentialsummen über wohlseparierte Punkte.
Ihre Hauptmerkmale sind:

\begin{enumerate}
  \item \textbf{Schärfe:} Die Konstante $(N + Q^2)$ ist bestmöglich.
  \item \textbf{Flexibilität:} Sie gilt für beliebige komplexe Folgen $(a_n)$
    ohne Strukturvoraussetzungen.
  \item \textbf{Charakter-Version:} Über Gauß-Summen und Farey-Brüche
    überträgt sich die Ungleichung direkt auf Dirichlet-Charaktersummen.
  \item \textbf{Konsequenzen:} Der Satz von Bombieri--Vinogradov (eine
    gemittelte Form der GRH), Linniks Satz über die kleinste Primzahl in
    einer Progression und Schranken für Charaktersummenmittel folgen alle
    daraus.
\end{enumerate}

Das Große Sieb bleibt eines der zentralen Werkzeuge der modernen analytischen
Zahlentheorie und findet weiterhin neue Anwendungen --- von der additiven
Kombinatorik (Bogolyubov--Ruzsa) bis zur Zufallsmatrizentheorie.

\begin{thebibliography}{10}

\bibitem{Bombieri1965}
E.~Bombieri.
\newblock On the large sieve.
\newblock \emph{Mathematika}, 12:201--225, 1965.

\bibitem{MV1973}
H.~L.~Montgomery und R.~C.~Vaughan.
\newblock The large sieve.
\newblock \emph{Mathematika}, 20:119--134, 1973.

\bibitem{IwaniecKowalski2004}
H.~Iwaniec und E.~Kowalski.
\newblock \emph{Analytic Number Theory}.
\newblock American Mathematical Society, Providence, RI, 2004.

\bibitem{CojocMurty2005}
A.~C.~Cojocaru und M.~R.~Murty.
\newblock \emph{An Introduction to Sieve Methods and Their Applications}.
\newblock Cambridge University Press, Cambridge, 2005.

\bibitem{Davenport2000}
H.~Davenport.
\newblock \emph{Multiplicative Number Theory}, 3.~Aufl.
\newblock Springer, New York, 2000.

\bibitem{Montgomery1971}
H.~L.~Montgomery.
\newblock \emph{Topics in Multiplicative Number Theory}.
\newblock Lecture Notes in Mathematics~227. Springer, Berlin, 1971.

\bibitem{Linnik1944}
Yu.~V.~Linnik.
\newblock On the least prime in an arithmetic progression.
  I.~The basic theorem.
\newblock \emph{Rec.\ Math.\ [Mat.\ Sb.] N.S.}, 15(57):139--178, 1944.

% BUG-P15-004 behoben (DE): Zhang2013 → Zhang2014 (Druckjahr 2014).
\bibitem{Zhang2014}
Y.~Zhang.
\newblock Bounded gaps between primes.
\newblock \emph{Ann.\ of Math.}, 179(3):1121--1174, 2014.

% BUG-P15-005 behoben (DE): Xylouris (2011) als Beleg für L ≤ 5 hinzugefügt.
\bibitem{Xylouris2011}
T.~Xylouris.
\newblock On the least prime in an arithmetic progression and estimates for
  the zeros of Dirichlet $L$-functions.
\newblock \emph{Acta Arith.}, 150:65--91, 2011.

\end{thebibliography}

\end{document}
